\section{Tensor Products of Representations}

  \begin{definition}[Tensor Product of Representations]
    \label{def:tensor-product-of-representations}
    If $V$ and $W$ are two different representations of a group $\mathcal{G}$, then we know that $V \oplus W$ is also a representation of $\mathcal{G}$. Furthermore, the tensor product space $V \otimes W$ also defines a representation of $\mathcal{G}$. That is, given representations
    \begin{align*}
      & \rho_V: \mathcal{G} \rightarrow \GL(V) \\
      & \rho_W: \mathcal{G} \rightarrow \GL(W)
    \end{align*}
    The homomorphism $\rho_V \otimes \rho_W: \mathcal{G} \rightarrow \GL(V \otimes W)$ is also a representation of $\mathcal{G}$, which is defined
    \begin{equation}
      (\rho_V \otimes \rho_W)(g) (v \otimes w) \coloneqq \rho_V (g) (v) \otimes \rho_W (g) (w)
    \end{equation}
    or represented in shorthand notation, 
    \begin{equation}
      g(v \otimes w) \coloneqq (g v) \otimes (g w)
    \end{equation}
    We know that $\exp(H)$ acts on $V$ and $W$ since it is an element of $\GL(V)$ and $\GL(W)$. This means that
    \begin{equation}
      \exp(H)(v \otimes w) \coloneqq \big( \exp(H)(v)\big) \otimes \big( \exp(H)(w)\big)
    \end{equation}
    If $H$ ($= \rho_V (H)$ or $\rho_W(H)$) has an eigenvalue $\lambda$ on $v$ in $V$ and eigenvalue $\mu$ on $w$ in $W$, then 
    \begin{equation}
      \exp(H) (v \otimes w) = (e^\lambda v) \otimes (e^\mu w) = e^{\lambda + \mu} v \otimes w
    \end{equation}
    That is, eigenvalues of $H$ \textbf{add} on tensor products. 
  \end{definition}

  \begin{example}[Weights of the $2$-Dimensional Representation]
    \label{ex:weights-of-the-2-dimensional-representation}
    Recall that the $2$ dimensional representation $V$ of $\mathfrak{sl}_2 \mathbb{C}$ has eigenvalues $1$ and $-1$ (with corresponding eigenvectors $e_1$ and $e_{-1}$). So, $V \otimes V$ has eigenvalues 
    \begin{align*}
      & (-1) + (-1) = -2, \;\; (-1) + 1 = 0 \\
      & 1 + (-1) = 0, \;\; 1 + 1 = 2
    \end{align*}
    Therefore, the eigenvalues of $V \otimes V$ is $-2$ (geometric multiplicity of 1), $0$ (geometric multiplicity of 2), and $2$ (geometric multiplicity of 1), (Notation-wise, the $n$-dimensional \hyperref[def:irreducible-representation]{irreducible representation} of $\mathfrak{sl}_2 \mathbb{C}$ is denoted $\mathbf{n}$.) which means that
    \begin{equation}
      \mathbf{2} \otimes \mathbf{2} = \mathbf{3} \oplus \mathbf{1}
    \end{equation}
    We can decompose $V \otimes V$ into its symmetric and exterior power components. Sym$^2 V$ has basis (of eigenvectors)
    \begin{equation}
      \{e_{-1} \odot e_{-1}, \; e_{-1} \odot e_1, \; e_1 \odot e_1\}
    \end{equation}
    where the corresponding eigenvalues are $-2$, $0$, and $2$, respectively. So, $\dim{Sym^2 V} = 3$, which means that $Sym^2 V = \mathbf{3}$. As for the exterior power component of $V$, $\Lambda^2 V$ has basis $\{e_{-1} \wedge e_1\}$ with eigenvalue $= 0 \implies \dim{\Lambda^2 V} = 1$, meaning that $\Lambda^2 V = \mathbf{1}$. Therefore, 
    \begin{equation}
      V \otimes V = Sym^2 V \oplus \Lambda^2 V = \mathbf{3} \oplus \mathbf{1}
    \end{equation}
  \end{example}
